\subsection{带赋值的域的完备化}

命题 5. 设 K 是一个未必交换的域，v 是 K 上的一个赋值，G 是赋以离散拓扑的群 \(v(\mathbf{K}^{*})\) 。

\begin{enumerate}
\def\labelenumi{(\alph{enumi})}
\item
  完备化环 \(\hat{\mathbf{K}}\) （即 \(\mathbf{K}\) 赋以 \(\mathcal{T}_v\) 后的完备化）是一个拓扑域。
\item
  映射 \(v: \mathbf{K}^* \to \mathbf{G}\) 可以唯一地延拓为连续映射 \(\hat{v}: \hat{\mathbf{K}}^* \to \mathbf{G}\) 。映射 \(\hat{v}\) （用 \(\hat{v}(0) = +\infty\) 加以延拓）是 \(\hat{\mathbf{K}}\) 上的一个赋值，且 \(\hat{v}(\hat{\mathbf{K}}^*) = v(\mathbf{K}^*)\) 。
\item
  \(\hat{\mathbf{K}}\) 上的拓扑是由赋值 \(\hat{v}\) 定义的拓扑。
\item
  对一切 \(\alpha \in \mathbf{G}\) ，设 \(\mathrm{V}_{\alpha}, \mathrm{V}_{\alpha}'\) 是 \(\mathbf{K}\) 的子群，由条件 \(v(x) > \alpha\) 、\(v(x) \geqslant \alpha\) 定义。那么 \(\overline{\mathrm{V}}_{\alpha}, \overline{\mathrm{V}}_{\alpha}'\) （\(\mathrm{V}_{\alpha}, \mathrm{V}_{\alpha}'\) 在 \(\hat{\mathbf{K}}\) 中的闭包）分别由条件 \(\hat{v}(x) > \alpha\) 、\(\hat{v}(x) \geqslant \alpha\) 定义。
\item
  \(\hat{v}\) 的环是 \(\hat{\mathbf{A}}\) ，即环 \(\mathbf{A}\) （\(v\) 的环）的完备化；\(\hat{v}\) 的理想是 \(\hat{\mathfrak{m}}\) ，即理想 \(\mathfrak{m}\) （\(v\) 的理想）的完备化。
\item
  \(\hat{\mathbf{A}} = \mathbf{A} + \hat{\mathfrak{m}}\) ；\(\hat{v}\) 的剩余域典范地等同于 \(v\) 的剩余域。
\end{enumerate}

为证明 (a) ，只需（《一般拓扑学》第三章 § 6 第 8 节命题 7）证明下述结论：设 \(\mathfrak{F}\) 是 \(\mathbf{K}^*\) 上（关于加法一致结构的）一个 Cauchy 滤子，0 不是它的触点；那么 \(\mathfrak{F}\) 在双射 \(x \mapsto x^{-1}\) 下的像是一个（关于加法一致结构的）Cauchy 滤子。因为 0 不是 \(\mathfrak{F}\) 的触点，故存在 \(\mathbf{M} \in \mathfrak{F}\) 与 \(\beta \in \mathbf{G}\) 使得 \(\beta\) 是 \(v(\mathbf{M})\) 的一个上界。设 \(\alpha \in \mathbf{G}\) 。若 \(\mathbf{M}'\) 是 \(\mathfrak{F}\) 的一个含于 \(\mathbf{M}\) 的元素，且 \(v(x - y) > \sup (\alpha + 2\beta, \beta)\) 对 \(x \in \mathbf{M}'\) 与 \(y \in \mathbf{M}'\) 成立，则 \(v(x^{-1} - y^{-1}) > \alpha\) 对 \(x \in \mathbf{M}'\) 与 \(y \in \mathbf{M}'\) 成立（第 1 节引理 1）。由此得 (a) 。

由第 1 节命题 1，\(v|K^*\) 是 \(K^*\) 到 \(G\) 的连续同态，因此可以唯一地延拓为连续同态 \(\hat{v}\) （由 \(K^*\) 到 \(G\) ）。关系式

\[
\hat {v} (x + y) \geqslant \inf (\hat {v} (x), \hat {v} (y))
\]

在 \(\mathbf{K}^*\) 中成立，因而由连续性在 \(\hat{\mathbf{K}}^*\) 中也成立。于是 \(\vartheta\) （用 \(\hat{\upsilon}(0) = +\infty\) 加以延拓）是 \(\hat{\mathbf{K}}\) 上的一个赋值，(b) 得证。

现在证明 (d) 。设 \(\alpha \in G\) ，\(x \in \overline{V}_{\alpha} - \{0\}\) 。当 \(y\) 在 \(V_{\alpha}\) 中充分接近 \(x\) 时，有 \(\hat{v}(x) = \hat{v}(y) = v(y)\) ，因而 \(\hat{v}(x) > \alpha\) 。反之，设 \(x \in \widehat{K}^*\) 满足 \(\hat{v}(x) > \alpha\) ；当 \(y\) 在 \(K^*\) 中充分接近 \(x\) 时，有 \(v(y) = \hat{v}(y) = \hat{v}(x)\) ，因此 \(y \in V_{\alpha}\) ，于是 \(x \in \overline{V}_{\alpha}\) 。这样，\(\overline{V}_{\alpha}\) 是由这样的 \(x \in \widehat{K}\) 组成的集合：\(\hat{v}(x) > \alpha\) 。对 \(V_{\alpha}'\) 论证类似。这就证明了 (d) 。

考虑到《一般拓扑学》第三章 § 3 第 4 节命题 7，断言 (c) 是 (d) 的推论。断言 (e) 是 (d) 的特例。最后设 \(x \in \hat{\mathbf{A}}\) ；存在 \(y \in \mathbf{A}\) 使得 \(\hat{v}(x - y) > 0\) ；于是 \(z = x - y \in \hat{\mathfrak{m}}\) ，因而 \(x = y + z \in \mathbf{A} + \hat{\mathfrak{m}}\) ；这样 \(\hat{\mathbf{A}} = \mathbf{A} + \hat{\mathfrak{m}}\) ，这证明了 (f) 。

注. 对一切 \(x \in \hat{\mathbf{K}}\) （不属于 \(\hat{\mathbf{A}}\) ），存在 \(x_0 \in \mathbf{K}\) 使得 \(\hat{v}(x - x_0) > 0\) ，\(\hat{v}(x) = \hat{v}(x_0) = v(x_0) < 0\) ；于是 \(x_0^{-1}x \in \hat{\mathbf{A}}\) ，且由于 \(x_0^{-1} \in \mathbf{A}\) ，可见若令 \(S = A - \{0\}\) ，就可以写成 \(\hat{\mathbf{K}} = S^{-1}\hat{\mathbf{A}}\) 。

